% STAGE08_PUBLIC_INTERFACE_DEMO — template stress test.
\DocumentMetadata{tagging=on,tagging-setup={math/setup=tex-AF},pdfversion=2.0,pdfstandard=ua-2,lang=zh-CN}
\documentclass[UTF8,a4paper,11pt,openany]{ctexbook}
\usepackage{statlearnbook}
\SLSetBookMetadata{模板样章}{统计学习方法初学者讲义}{公共环境、编号与交叉引用验证}
\begin{document}
\frontmatter
\SLMakeTitle
\SLPrintContents
\mainmatter
\chapter{公共模板验证}\label{chap:template-demo}
\SLChapterOpening{从二次损失的定义进入梯度下降，再核对算法、图表和长内容的统一版式}{公共接口能否同时承载严格数学对象与稳定的读者导航}{识别章节结构、跟随一次梯度更新，并验证交叉引用与索引}{向量内积和一元函数求导}{若内积或求导自检失败，先回看先修框；其余对象按本章分级标题逐项复核}
\chaptergoals{验证中文字体、数学字体、定理环境、算法、图表、书签、交叉引用与索引。}
\begin{prerequisitebox}
读者只需熟悉向量内积和一元函数求导。
\end{prerequisitebox}
\begin{precheckbox}
请确认能够计算 $\inner{x,y}=x^{\mathsf T}y$，并说明 $\norm{x}_2^2=\inner{x,x}$。
\end{precheckbox}
\begin{dependencybox}
向量内积与一元微分 $\longrightarrow$ 二次损失 $\longrightarrow$ 梯度下降。
\end{dependencybox}
\begin{symboltablebox}
$y$：真实值；$\hat y$：预测值；$L$：损失函数；$w$：待更新参数。
\end{symboltablebox}
\SLTeachSection{定义、公式与定理}{sec:template-definition}
\knowledgeanchor{TPL-DEF-001}{二次损失}
\begin{definition}[二次损失]\label{def:quadratic-loss}
对预测值 $\hat y\in\R$ 和真实值 $y\in\R$，定义
\begin{equation}\label{eq:quadratic-loss}
  L(y,\hat y)=\frac{1}{2}(y-\hat y)^2.
\end{equation}
\end{definition}
\begin{theorem}[非负性]\label{thm:loss-nonnegative}
\cref{eq:quadratic-loss} 定义的损失满足 $L(y,\hat y)\ge 0$，且等号成立当且仅当 $y=\hat y$。
\end{theorem}
\begin{proof}
实数平方非负；平方为零当且仅当其底数为零。
\end{proof}
由\cref{def:quadratic-loss,thm:loss-nonnegative}可见，编号和交叉引用均按章组织。

\begin{geometrybox}
二次损失是误差到原点距离平方的一半。
\end{geometrybox}
\begin{probabilitybox}
在方差固定的高斯噪声模型中，最小二乘与最大似然等价。
\end{probabilitybox}
\begin{optimizationbox}
\termindex{二次损失}为凸函数，梯度为 $\hat y-y$。
\end{optimizationbox}

\SLDirectSection{算法、图与表}{sec:template-assets}
\begin{algorithm}[H]
\KwIn{初值 $w_0$、步长 $\eta>0$、迭代次数 $T$}
\KwOut{参数 $w_T$}
$w\leftarrow w_0$\;
\For{$t\leftarrow 0$ \KwTo $T-1$}{
  $w\leftarrow w-\eta\nabla f(w)$\;
}
\KwReturn{$w$}\;
\caption{梯度下降模板}\label{alg:gradient-template}
\end{algorithm}
\begin{warningbox}
若步长不是正数则拒绝执行；$T=0$ 时直接返回 $w_0$，用于覆盖合法退化输入。
\end{warningbox}
\begin{complexitybox}
训练（一次更新）为 $O(d)$ 时间；预测复杂度取决于后续模型，本更新模板不单独定义预测；参数存储为 $O(d)$ 空间。
\end{complexitybox}

\begin{figure}[H]
\centering
\begin{tikzpicture}[>=Stealth]
  \draw[->] (-0.2,0)--(4.2,0) node[right] {$x$};
  \draw[->] (0,-0.2)--(0,3.0) node[above] {$f(x)$};
  \draw[SLDeepBlue,thick,domain=0.2:3.8,samples=80] plot (\x,{0.22*(\x-2)^2+0.35});
  \fill[SLMidBlue] (2,0.35) circle (2pt) node[below=4pt] {最小点};
\end{tikzpicture}
\caption{克制配色的可复现示意图}\label{fig:template-curve}
\end{figure}
\cref{fig:template-curve}直观显示了凸二次函数的唯一最小点。

\begin{table}[H]
\centering
\caption{环境编号测试}\label{tab:template-environments}
\begin{tabular}{lll}
\toprule
对象 & 编号方式 & 引用示例\\
\midrule
公式 & 按章 & \cref{eq:quadratic-loss}\\
算法 & 按章 & \cref{alg:gradient-template}\\
图 & 按章 & \cref{fig:template-curve}\\
\bottomrule
\end{tabular}
\end{table}

\SLDirectSection{例题、解释与自检}{sec:template-example}
\begin{example}\label{ex:template-number}
当 $y=2,\hat y=1$ 时，计算二次损失。
\end{example}
\begin{analysisnote}
先代入\cref{eq:quadratic-loss}，再计算差值的平方。
\end{analysisnote}
\begin{solution}
$L(2,1)=\tfrac12(2-1)^2=\tfrac12$。
\end{solution}
\begin{warningbox}
不要漏掉定义中的系数 $\tfrac12$；它用于简化求导结果。
\end{warningbox}
\SLAdvancedSection{分页环境压力测试}{sec:template-pagination}
\begin{derivationgoalbox}
验证较长推导准备、长表格与章末结构环境在页边界处保持完整。
\end{derivationgoalbox}
\begin{derivationprepbox}
已知 $w\in\R^d$、$\eta>0$ 与可微函数 $f:\R^d\to\R$。下列对象只用于模板排版测试，不作为教材数学结论。
\end{derivationprepbox}
\begin{longtable}{clp{0.55\textwidth}}
\caption{长表格分页与重复表头测试}\label{tab:long-template}\\
\toprule
序号 & 对象 & 排版检查\\
\midrule
\endfirsthead
\multicolumn{3}{c}{\small 表\thetable\ （续）}\\
\toprule
序号 & 对象 & 排版检查\\
\midrule
\endhead
\bottomrule
\endfoot
1 & 定义 & 标题、编号、浅底与左侧细线\\
2 & 定理 & 条件、结论与证明衔接\\
3 & 公式 & 按章编号并保留足够行宽\\
4 & 算法 & 输入、输出、循环与返回值\\
5 & 图 & 坐标轴、曲线、标签与图注\\
6 & 表 & 表头、对齐、边界与续表标记\\
7 & 例题 & 题面、思路、计算与答案\\
8 & 练习 & 题目来源类别与完整解析\\
9 & 符号 & 含义、维数、取值空间与首次定义\\
10 & 引用 & 中文对象名、编号与可点击链接\\
11 & 页眉 & 册名、章名、节名与页码\\
12 & 索引 & 符号索引与主题索引\\
13 & 书签 & 章、节层级与中文标题\\
14 & 长内容 & 允许在页边界处自然续排\\
15 & 留白 & 不产生异常整页空白\\
16 & 字体 & 宋体、黑体、Times与协调数学字体\\
17 & 颜色 & 黑色正文与克制深色标题\\
18 & 边距 & A4页面与合理版心\\
\end{longtable}
\begin{chapterexercisebox}
\SLSourceBookExercises
\textnormal{模板说明：此处只验证“原书练习整理”分类的版式；样题并非原书题，正式章节必须依据专门索引核准来源。}
\begin{exercise}\label{exr:template-gradient}
求 $L(y,\hat y)$ 关于 $\hat y$ 的导数。
\end{exercise}
\SLLectureSupplementExercises
\begin{exercise}\label{exr:template-risk}
指出\cref{tab:long-template}中覆盖的四类高风险排版对象。
\end{exercise}
\end{chapterexercisebox}
\begin{chapteranswerbox}
\SLExerciseSolutionHeading{exr:template-gradient}
\begin{solution}
$\partial L/\partial\hat y=\hat y-y$。
\end{solution}
\SLExerciseSolutionHeading{exr:template-risk}
\begin{solution}
公式、算法、图和长表格均已在样章中实际排版；验收时还需同时检查编译日志与页面图像。
\end{solution}
\end{chapteranswerbox}
% Natural-language example for a chapter without executable steps.
\begin{notapplicablebox}[title=契约示例：确实不适用]
\SLNotApplicableItem{算法说明}{若某章只约定数学符号且没有可执行过程，可直接说明该章无需给出算法步骤；本例只演示排版，不代表任何正文结论。}
\end{notapplicablebox}
\SLChapterFiveSummary{二次损失与公共排版接口}{由平方非负性得到损失非负性，并由求导得到梯度更新}{按输入、循环、返回值执行梯度下降模板}{需要验证教材环境、编号、图表、分页、目录与索引时}{回到\cref{sec:template-definition,sec:template-assets}逐项核对，并用章末自检确认结果}
\begin{selfcheckbox}
能够解释\cref{eq:quadratic-loss}的非负性，并手工跟随\cref{alg:gradient-template}的一次更新。
\end{selfcheckbox}
\index[symbols]{L@$L$：损失函数}
\backmatter
\printbookindexes
\end{document}
